% ============================================================
% fig_hopfield_ejemplo.tex: Ejemplo numérico paso a paso (4 neuronas)
% ============================================================
\begin{figure}[H]
\centering
\begin{tikzpicture}[
    cell/.style={draw, minimum width=0.55cm, minimum height=0.55cm, font=\footnotesize},
    on/.style={cell, fill=blue!60, text=white},
    off/.style={cell, fill=gray!20},
    lbl/.style={font=\footnotesize\bfseries},
    stepnum/.style={font=\small\bfseries, text=blue!70},
    arrow/.style={-{Stealth}, thick, gray}
]

% ── Fila 0: encabezados ─────────────────────────────────────
\node[lbl] at (-1.8, 3.8) {Paso};
\node[lbl] at ( 0.4, 3.8) {Estado $\bm{s}$};
\node[lbl] at ( 3.6, 3.8) {Descripción};

% ── Paso 0: patrón almacenado ───────────────────────────────
\node[stepnum]  at (-1.8, 3.0) {0};
\foreach \v/\i in {+1/0, -1/1, +1/2, -1/3}{
    \pgfmathparse{(\v>0)?1:0}
    \ifnum\pgfmathresult=1
        \node[on]  at (\i*0.65, 3.0) {\v};
    \else
        \node[off] at (\i*0.65, 3.0) {\v};
    \fi
}
\node[font=\footnotesize, align=left] at (5.9, 3.0)
    {Patrón almacenado $\xi^1 = [+1,-1,+1,-1]$};

% ── Separador ───────────────────────────────────────────────
\draw[dashed, gray!50] (-2.5, 2.55) -- (5.0, 2.55);

% ── Paso 1: entrada ruidosa ─────────────────────────────────
\node[stepnum]  at (-1.8, 2.1) {1};
\foreach \v/\i in {+1/0, +1/1, +1/2, -1/3}{
    \pgfmathparse{(\v>0)?1:0}
    \ifnum\pgfmathresult=1
        \node[on]  at (\i*0.65, 2.1) {\v};
    \else
        \node[off] at (\i*0.65, 2.1) {\v};
    \fi
}
\draw[red!70, very thick] (0.65-0.27, 2.1-0.27) rectangle (0.65+0.27, 2.1+0.27);
\node[font=\footnotesize, align=left] at (5.6, 2.1)
    {Entrada ruidosa: $s_2$ invertido {\color{red!70}$\leftarrow$ error}};

% ── Flecha de actualización ─────────────────────────────────
\draw[arrow] (0.95, 1.78) -- (0.95, 1.35);
\node[font=\tiny, gray] at (2.5, 1.57) {actualización: $s_i \leftarrow \text{sgn}\!\left(\sum_j w_{ij} s_j\right)$};

% ── Paso 2: después de una iteración ────────────────────────
\node[stepnum]  at (-1.8, 1.05) {2};
\foreach \v/\i in {+1/0, -1/1, +1/2, -1/3}{
    \pgfmathparse{(\v>0)?1:0}
    \ifnum\pgfmathresult=1
        \node[on]  at (\i*0.65, 1.05) {\v};
    \else
        \node[off] at (\i*0.65, 1.05) {\v};
    \fi
}
\node[font=\footnotesize, align=left] at (6.0, 1.05)
    {Tras actualizar $s_2$: patrón recuperado $\xi^1$ {\color{green!60!black}$\checkmark$}};

% ── Estado estable ───────────────────────────────────────────
\draw[dashed, gray!50] (-2.5, 0.6) -- (6.8, 0.6);
\node[font=\footnotesize, gray, align=left] at (2.0, 0.25)
    {Estado estable: $E$ alcanzó mínimo local $\Rightarrow$ la red converge.};

% ── Leyenda ─────────────────────────────────────────────────
\node[on,  label=right:{\footnotesize $s_i = +1$ (activa)}]  at (-1.2, -0.35) {};
\node[off, label=right:{\footnotesize $s_i = -1$ (inactiva)}] at ( 2.8, -0.35) {};

\end{tikzpicture}
\caption{Ejemplo de recuperación de un patrón en una red de Hopfield con
4 neuronas. El patrón almacenado $\xi^1 = [+1,-1,+1,-1]$ se recupera
correctamente a partir de una versión ruidosa en la que el bit $s_2$ estaba
invertido. La regla de actualización $s_i \leftarrow \text{sgn}\!\left(\sum_j w_{ij} s_j\right)$
corrige el error en una sola iteración.}
\label{fig:hopfield_ejemplo}
\end{figure}
